% From chips or fibres: hot pressing of particle board vs continuous-fibre material extrusion
\documentclass[tikz,border=6pt]{standalone}
\usepackage{xfab-tikz}
\begin{document}
\begin{tikzpicture}[x=1cm,y=1cm]
\xfabpanels{Hot pressing of boards}{a mat of resin-coated chips is pressed and cured}
           {Continuous-fibre extrusion}{a fibre tow is laid inside each polymer road}

% ================= LEFT: particle board press
\begin{scope}[shift={(0.5,-5.6)}]
  \draw[machinedark] (0,3.0) rectangle (5.8,3.5);
  \draw[machinedark] (0,0.0) rectangle (5.8,0.5);
  \foreach \xx in {0.6,1.6,...,5.6}{\fill[heat] (\xx-0.25,2.92) rectangle (\xx+0.25,3.0); \fill[heat] (\xx-0.25,0.5) rectangle (\xx+0.25,0.58);}
  % mat of chips
  \fill[fibrous!10] (0.3,0.58) rectangle (5.5,2.0);
  \pgfmathsetseed{21}
  \foreach \i in {1,...,90}{
    \pgfmathsetmacro{\fx}{0.4+4.8*rnd}\pgfmathsetmacro{\fy}{0.7+1.15*rnd}\pgfmathsetmacro{\fa}{-25+50*rnd}
    \draw[fibrous!70!black, line width=1.1pt] (\fx,\fy) -- ++(\fa:0.35);}
  % upper platen moves down
  \draw[flow=steel] (4.6,4.2) -- (4.6,3.6);
  \draw[flow=steel] (4.6,-0.7) -- (4.6,-0.1);
  \node[lbl, anchor=west] at (4.8,4.0) {pressure};
  \node[lbl, anchor=west] at (0.0,3.75) {heated platens ($\approx$ 200\,°C)};
  \node[lbl, anchor=west] at (0.3,2.3) {wood chips + resin};
  \node[lbl, anchor=north west] at (0.3,-0.2) {resin cures under heat and pressure};
\end{scope}

% ================= RIGHT: continuous fibre co-extrusion
\begin{scope}[shift={(8.6,-6.1)}]
  \draw[machine] (0,0) rectangle (5.6,0.25);
  % roads in cross-section with a fibre bundle in each
  \foreach \l in {0,...,3}{
    \foreach \b in {0,...,4}{
      \pgfmathsetmacro{\cx}{0.9+\b*0.56}\pgfmathsetmacro{\cy}{0.47+\l*0.44}
      \fill[plastic!35, draw=plastic!70!black, line width=0.4pt] (\cx,\cy) ellipse (0.28 and 0.22);
      \foreach \dx/\dy in {-0.07/0.03, 0.04/0.05, 0.0/-0.05, 0.09/-0.02, -0.08/-0.06}
        \fill[ink] (\cx+\dx,\cy+\dy) circle (0.022);}}
  % nozzle laying the next road (side view, above stack)
  \draw[machine] (3.5,2.7) rectangle (4.7,3.3);
  \fill[energy!80] (3.65,2.7) -- (4.55,2.7) -- (4.18,2.15) -- (4.02,2.15) -- cycle;
  \fill[plastic!35, draw=plastic!70!black, line width=0.4pt, rounded corners=3pt] (1.2,2.0) rectangle (4.15,2.28);
  \draw[ink, line width=1pt] (1.2,2.14) -- (4.1,2.14) -- (4.1,2.6);
  % feeds
  \draw[plastic!60, line width=3pt] (3.85,5.0) -- (3.85,3.3);
  \draw[ink, line width=1.2pt] (4.35,5.0) .. controls (4.35,4.3) .. (4.15,3.3);
  \draw[machinedark] (4.85,5.0) circle (0.35); \fill[white] (4.85,5.0) circle (0.1);
  \draw[ink, line width=1.2pt] (4.55,4.85) -- (4.35,4.7);
  \draw[dblmotion] (1.3,2.75) -- (2.6,2.75); \node[lbl] at (1.95,2.97) {$x,y$};
  % cutter
  \node[lbl, anchor=east] at (3.75,4.55) {polymer filament};
  \node[lbl, anchor=west] at (5.25,5.0) {fibre tow};
  \node[lbl, anchor=west] at (5.25,4.65) {(carbon, glass)};
  \node[lbl, anchor=west] at (4.8,1.3) {cross-section:};
  \node[lbl, anchor=west] at (4.8,1.0) {fibres in};
  \node[lbl, anchor=west] at (4.8,0.7) {every road};
  \node[lbl, anchor=north] at (2.8,-0.05) {fibres follow the tool path: strong along it, weak across};
\end{scope}
\xfabstate{fibrous}{1.5}{chip or fibre}
\end{tikzpicture}
\end{document}
